\section{Summary}\label{sec:summary}

A coding agent such as Amplifier runs on a \term{host model}: the large model that answers each turn of a
conversation. The fast-decisions bundle can \term{route} some of the work to a cheaper model (here Sonnet 5)
when it judges the work easy enough. Routing looks like an obvious saving. It is not obvious once a
conversation gets long, because of \term{prompt caching}: the provider keeps a cheap-to-reuse copy of the
conversation \emph{for each model separately}, so switching models means paying again to cache the whole
history (\cref{sec:caching}).

This report answers one question with measurements rather than estimates: \emph{what does routing an agent
to a cheaper model actually cost, in money and in quality, over realistic multi-turn sessions?}

\begin{keyidea}
\textbf{What we found} (test split, preregistered hypotheses; prices of \cref{tab:prices};
\cref{sec:confirm}):
\begin{itemize}[leftmargin=1.2em]
  \item \textbf{Fable 5.1 host: Sonnet is cheaper.} Choosing the model once at the start of the session
    (\term{sticky}) cost \CfFableStickyPairTwo\X\ the plain host (95\,\% CI \CfFableStickyPairCITwo), a
    saving of about \CfFableStickySaving\,\%, with turn-pass fraction inside the preregistered margin on
    the test scenarios. The shipped router, which re-chooses its model at the start of each turn, cost
    \CfFableShippedPairTwo\X; plain Sonnet for the whole session cost \CfFableSonnetPairTwo\X.
  \item \textbf{Opus 5.5 host: routing to Sonnet cost more.} Sticky \CfOpusStickyPairTwo\X, shipped
    \CfOpusShippedPairTwo\X, plain Sonnet \CfOpusSonnetPairTwo\X\ (pair reading); without the quality filter
    (arm reading) \CfOpusStickyArmTwo\X, \CfOpusShippedArmTwo\X\ and \CfOpusSonnetArmTwo\X, that is
    \CfOpusStickyArmExtra--\CfOpusSonnetArmExtra\,\% more. This depends on Opus's cache-read price: at about
    \$\BEStickyPair\ per million tokens instead of \$\PriceOpusRead, sticky would break even
    (\cref{sec:breakeven}).
  \item \textbf{Sticky cost less than shipped on both hosts:} \HThreeFable\X\ on Fable and \HThreeOpus\X\ on
    Opus. The preregistered test was only that sticky costs no more than shipped; on Fable the margin is small
    (upper bound \HThreeFableUpper).
  \item \textbf{The savings model passes a pooled check.} Fitted on \NTrain\ training scenarios, it
    predicted the held-out total within its 90\,\% interval and its intervals covered the held-out pairs at
    the preregistered rate. Its per-host and per-task figures are not reliable (\cref{sec:modellimits}).
\end{itemize}
\end{keyidea}

All \NHypConfirmed\ preregistered hypotheses were confirmed, under both readings of ``non-inferior''
quality that the preregistration allows (\cref{sec:confirm}). Three cautions travel with the result.
\begin{itemize}
  \item \textbf{Quality is close to the line.} Sticky on Fable passed the 5-point turn-pass margin on the test
    scenarios by \QFableStickyMarginGap; across all \QAllScen\ scenarios the margin was not shown
    (\QAllFableSticky, 95\,\% CI \QAllFableStickyCI), and routed Fable arms passed the final hidden tests
    \FinalGapLow--\FinalGapHigh\ points less often (exploratory; \cref{sec:finalpass}).
  \item \textbf{The router adds little beyond ``use Sonnet''.} \NStickySonnetOnly\ of \NStickySessions\ sticky
    sessions ran Sonnet only, and plain Sonnet cost about the same on Fable (\cref{sec:recs}).
  \item \textbf{Scope.} Sessions are scripted, at most \MaxTurns\ turns long, on one provider and one price
    table (\cref{sec:limits}).
\end{itemize}

\paragraph{How to read this report.} \Cref{sec:why} explains why estimates were not enough, and
\cref{sec:caching} explains prompt caching with a small diagram and a direct probe of the provider.
\Cref{sec:design} describes the experiment and \cref{sec:running} how it was run safely. \Cref{sec:confirm}
gives the confirmatory results; \cref{sec:explore} the exploratory ones, labelled as such.
\Cref{sec:recs,sec:limits,sec:repro} give recommendations, limitations and how to reproduce every number.
A glossary is in \cref{app:glossary}.
